% ============================================================
% DRAFT FOR DAMIAN'S REVIEW — the rewritten "Scoring documents at
% national scale" section for Chapter 2, from the approved skeleton
% (thesis/ch2_scoring_section_outline.md), in the voice specified by
% thesis/WRITING_STYLE.md. On approval this replaces the current
% sec:p1_text_method material; numeric literals then convert to
% make_numbers macros before integration.
% ============================================================

\section{Scoring Documents at National Scale}
\label{sec:p1_scoring}

The gold standard scores described above exist for the 103 schools that
received a visit. Extending the measurement of school culture to the full
national population of around 3,330 secondary schools requires instruments
that can read the public record instead: inspection reports, behaviour policy
documents, school websites, and the interview transcripts themselves. This
section describes how those instruments were built, how they were tested, and
what was learned in the process. There are seven parts. The first explains
what a scoring instrument is and the two forms it can take. The second
explains who produced the reference labels the instruments were tuned
against, which is a point on which this thesis owes the reader complete
transparency. The third describes how the written scoring rules were
themselves tested before any scoring took place. The fourth sets out the
three tests every instrument had to face, and the fifth summarises what
those tests taught, several findings of which are, to my knowledge, new. The
sixth part describes the instruments that were finally adopted, and the
seventh deals with reproducibility. The full scoring rules for every
instrument are reproduced in \cref{app:2A:prompts}.

% ------------------------------------------------------------
\subsection{What a scoring instrument is, and the two ways to build one}
\label{sec:p1_scoring_arch}

Each instrument takes one document about one school and returns a score from
1 to 5 for a single construct: warmth, strictness, or teaching. All scoring
is performed by a large language model (LLM), reading the same document a
human rater would read. There are two quite different ways to set this up,
and the difference between them turns out to matter a great deal.

The first approach is a \textit{prose rubric}. The model is given a written
marking scheme of the kind an examiner would use, describing what a document
at each band looks like, and it is asked to read the document and return a
band along with its reasoning. The judgement is the model's own. The second
approach is a \textit{decomposed instrument}. Here the model is never asked
for a judgement at all. Instead it answers a series of narrow factual
questions about the document. For example, does the website name a
consequence for lateness? Can it list the sentences in which an adult and a
pupil appear together with something passing between them? The answers to
these questions are then converted into a band by simple arithmetic in
ordinary code, written and fixed in advance. In the first approach the model
is the examiner; in the second it is closer to a careful clerk, and the
examining is done by a formula that anyone can inspect.

It would be natural to assume that one of these designs is simply better
than the other. That assumption was tested extensively, and the answer is
that each wins in different places and for identifiable reasons, which are
set out in \cref{sec:p1_scoring_lessons} below.

% ------------------------------------------------------------
\subsection{Who produced the labels: model raters and a human anchor}
\label{sec:p1_scoring_raters}

Tuning an instrument requires reference labels: documents whose correct band
is already known, so that the instrument's answers can be checked and its
rules improved. It is important to be clear about where those labels came
from, because they were not, for the most part, produced by human raters.

The labelling was carried out by three independent instances of a larger
language model (Anthropic's Claude, Opus class), each given only the written
scoring rule and the documents. The three raters worked blind: none was told
that the others existed, none saw my labels or any model score, and each
received the documents in a different shuffled order. The reference label
for each document was the majority vote of the three. The scoring model
throughout, by contrast, was a much smaller and cheaper model (gpt-4o-mini),
chosen because reading several thousand documents per instrument has to be
affordable. The tuning exercise should therefore be understood precisely:
it asked whether a small model could be brought to reproduce a larger
model's reading of a written rule. Where this chapter reports agreement
between an instrument and its labels, that is agreement between two
artificial readers, and the tables say so.

My own role in the labelling was threefold. I wrote the scoring rules and
revised them; I made the boundary rulings that the raters could not, which
are recorded in an adjudication log; and I labelled the earlier development
packs myself, before the model-rater design was adopted. Those earlier
labels also provide a calibration that is worth stating: when I relabelled
the same documents two weeks apart, blind to my first pass, I agreed with
myself at a weighted kappa of $+0.770$. No rater, human or artificial, can
be expected to agree with me more closely than I agree with myself, so this
figure sets a ceiling of roughly $0.88$ on any agreement measured against my
labels, and it is the standard against which the agreement figures below
should be read.

Two safeguards make this design defensible rather than merely convenient.
The first is that the ultimate test of every instrument is never its
agreement with any rater, artificial or human, but its relationship to the
independently observed visit scores of the gold standard, and to external
records where these exist. Those criteria share nothing with either model
family. The second is that the risk of artificial raters drifting away from
the researcher's own judgement was checked rather than assumed. On the packs
where both my labels and the model raters' labels exist, agreement was
substantial; but in one measured incident a model adjudicator, asked to
resolve a disagreement, sided with the scoring model against me. That
incident is reported here deliberately. It shows that a rater from the same
model family as the scorer is not independent human ground truth, and it is
the reason the criterion tests in \cref{sec:p1_scoring_tests} are given the
final word throughout this chapter.

% ------------------------------------------------------------
\subsection{Testing the written rule before spending anything}
\label{sec:p1_scoring_rules}

A scoring rule that lives only in its author's head is not an instrument.
Before any model was tuned, each written rule was therefore tested for what
might be called determinacy: whether three cold readers, given nothing but
the rule and the documents, arrive at the same answers. Across six rule
packs the three blind raters agreed with each other at weighted kappas of
between $+0.77$ and $+0.96$. A rule that transfers this well is doing its
job; the person who wrote it is no longer needed for it to be applied.

Just as useful as the overall figure is what the disagreements revealed.
Wherever the three raters split on a document, the rung they split over was
by definition under-specified, and it was repaired before any scoring money
was spent. Two examples give the flavour. The rule for teaching quality in
inspection reports originally awarded its fourth band for a sentence pattern
("teachers check understanding and address misconceptions") that turned out
to be standard inspection boilerplate appearing in nearly every post-2022
report, so the raters were forced to award band 4 almost everywhere; the
band was re-cut so that the boilerplate sentence no longer reached it. The
rule for website strictness originally merged into one band the naming of a
virtue ("we value respect") and the naming of a conduct ("we expect
punctuality"), which the raters could not apply consistently; the band was
split along exactly that line. The principle that emerged is worth stating
plainly: a rule that needs its rater to invent a resolution is
under-specified, and the agreement it produces is agreement about the
inventions rather than the rule.

% ------------------------------------------------------------
\subsection{The three tests}
\label{sec:p1_scoring_tests}

Every instrument had to pass through three tests, in a fixed order, and the
order matters because each test can fail for a reason the previous one
cannot detect.

The first test is \textit{agreement}: does the instrument reproduce the
reference labels on documents it was not tuned on? The working bar was a
weighted kappa of at least $+0.70$ together with exact agreement of at least
55 per cent, with an additional guard against instruments that achieve
agreement by collapsing onto one or two bands. Tuning always took place on
one part of a labelled pack with the quoted figure computed on the untouched
remainder.

The second test is \textit{out-of-sample honesty}. An agreement figure
measured on the documents an instrument was tuned on always reads high, for
the same reason that a student who has seen the exam paper scores well.
Late in the project, every archived version of every instrument --- roughly
85 versions across eleven instrument families --- was scored once on freshly
labelled documents that no version had ever been fitted to. Every in-sample
figure had indeed read high, and several instruments that appeared to have
improved across their version histories had not improved at all. All
agreement figures reported in this chapter are from this out-of-sample
measurement.

The third test is \textit{validity}: does the score track something real
that is not a language model's opinion? For most constructs the criterion is
the gold standard itself, the 103 independently observed visits. For one
construct there is something better. The prominence of religious faith on a
school's website can be checked against the official register of religious
character, and the website instrument recovers the register almost
perfectly, with an area under the curve of 0.960. That result matters
beyond its own construct: it shows that when the signal is genuinely present
in the text, this method of reading it works. Where a construct fails
against its criterion, the failure is therefore informative about the text
rather than the machinery.

% ------------------------------------------------------------
\subsection{What the tests taught}
\label{sec:p1_scoring_lessons}

Five findings from this testing programme deserve to be reported as results
in their own right, because each one changed how the instruments were built,
and because I am not aware of prior literature that documents them.

First, and most consequentially, \textbf{agreement and validity are
different axes, and improving one can spend the other}. The version of the
Ofsted strictness instrument that agreed best with the reference labels
tracked the independently observed visit scores at only $+0.07$; the prose
version, whose label agreement was lower, tracked them at $+0.43$. This
pattern was not an isolated case: across the eleven families the
best-agreeing version was usually not the most valid one. The practical
lesson is that an instrument tuned solely to match labels can quietly become
an excellent predictor of the labelling process and a poor measure of the
world. The instruments adopted in this chapter were chosen on all three
tests together, and this chapter never treats a high kappa as evidence of
validity.

Second, \textbf{ask the model for facts, not judgements}. Wherever a
question put to the model required an evaluation ("is behaviour management
a strength of this school?") the answer tracked the model's own generosity
more than the document. Wherever the question was factual ("does the policy
name a tracking system in which a pupil's position determines the next
sanction?") the answer tracked the document. The one instrument family
whose version history ran consistently uphill was the one whose decisive
revision replaced its judgement questions with factual ones.

Third, \textbf{a scoring rule must not rest on the genre's own
boilerplate}. In the national website corpus, 99.1 per cent of schools use
a care word somewhere in their text, so a warmth band that fires on caring
language is a band that contains almost every school in England. The same
applies to the standard phrases of inspection prose. Several early
instruments failed exactly this way, and every adopted rule was checked
against the national base rates of the phrases it relies on.

Fourth, \textbf{an instrument can manufacture statistical independence that
reality does not have}. In the visits, warmth and strictness correlate at
$+0.589$: warm schools tend also to be orderly ones. The prose instrument,
scoring both constructs in one reading of the document, reproduces this
correlation almost exactly ($+0.578$). The decomposed instruments, scoring
each construct in a separate call, return correlations between $+0.01$ and
$+0.16$. A reader of those instruments' output would conclude that warmth
and strictness are unrelated dimensions, and would be wrong. Any study that
scores related constructs with separate instruments and then interprets the
correlation between them should treat this as a warning.

Fifth, \textbf{the model is part of the instrument}. Moving one rubric,
unchanged, from one scoring model to another collapsed its label agreement
from $+0.63$ to $+0.12$. A rubric therefore has no performance of its own;
only a rubric-and-model pair does. This is why every score in this project
records the exact model version that produced it, and why one instrument
(faith prominence) is deliberately kept on an older model: its validation
belongs to that pairing, and re-scoring with a newer model would create a
new, unvalidated instrument, not an updated one.

% ------------------------------------------------------------
\subsection{The instruments adopted}
\label{sec:p1_scoring_adopted}

\Cref{tab:instruments_adopted} summarises the instrument finally adopted for
each source and construct, with its architecture, its out-of-sample
agreement, its validity result, and the caveat that must accompany its use.
Three points from the table should be drawn out.

For the inspection reports, all three constructs are scored by a single
prose rubric in one reading of a grade-stripped report body. This is the
only document instrument in the project with demonstrated validity against
the visits, and the convergence result it supports is reported in
\cref{sec:p1_text_valid}. For the behaviour policies, a prose rubric is
likewise used; it is the instrument with the strongest label agreement in
the project ($+0.757$), and its complete lack of validity against enacted
culture is not a defect but the chapter's policy-as-aspiration finding.
For the websites, decomposed instruments are used, and the warmth
instrument deserves a sentence of history: its floor question ("is there a
relational statement in this text?") was answered too conservatively by the
scoring model as a yes/no question, placing 78 per cent of schools in the
bottom band, and the repair was to ask the model to list and count the
qualifying sentences instead, with every listed sentence verified verbatim
against the document before being counted. Models, it turns out, gate
badly and count well. The final website strictness instrument carries five
bands for consistency with the other scales; its fifth band is defined by
two factual criteria and is genuinely rare, firing for 25 schools in
England, and no labelled pack contains enough such schools to verify the
band against raters. That caveat travels with any use of it.

One rule governs how these instruments may be used in \cref{ch:paper2} and
beyond: within any comparison, the two sides must share an architecture.
A strictness-versus-warmth contrast drawn inside one source must come from
instruments built the same way, because finding four above shows that
mixing architectures inside a contrast can manufacture or destroy the very
relationship being studied.

% ------------------------------------------------------------
\subsection{Reproducibility}
\label{sec:p1_scoring_repro}

Nothing in this pipeline is perfectly deterministic, and it is better to
measure the imperfection than to assert its absence. The scoring model's
test-retest reliability, measured by re-scoring identical documents, is
approximately 92 per cent exact agreement per construct. The older model
retained for the faith prominence instrument is considerably worse, at
roughly two identical scores in three, which is disclosed wherever that
instrument's output is used. Every score row in the dataset records the
prompt version, a digest of the exact rubric text, and the model version
that produced it, so any figure in this thesis can be traced to the
instrument that generated it. All printed numbers are generated from the
dataset by a script that also checks the chapters for superseded values, so
a stale number cannot silently survive a revision. Finally, one behaviour
policy document in the national corpus could not be scored at all: its PDF
carries a broken font encoding that yields unreadable text, and it is
recorded as unscoreable rather than being approximated.

% ============================================================
% Integration notes (not for the tex):
% - tab:instruments_adopted to be built from the OOS table + caveats.
% - Numeric literals convert to make_numbers macros at integration.
% - The exact Opus model identifier to be pinned from session records.
% ============================================================
